% conclusion.tex

The electric sector of the A=1 induced gauge action is now as concrete as the
magnetic one. Reading the on-site advance $\delta\phi=\omega+V(x)$ as a temporal
link turns the apparent obstruction---``no electric Wilson loop''---into a
temporal plaquette whose holonomy delivers the permittivity
$\varepsilon=P=\sum_a V_a V_a^{\!\top}$, eigenvalues $\{1,4,4\}$, optical axis
$(1,1,-1)$ suppressed: the exact mirror of the magnetic $\{4,4,16\}$. The block is
not only derived but read off the engine's real tick evolution
(\autoref{subsec:electric_block}), so it rests on the coupling the framework
actually implements, not an analogy. The two sectors are bound by an exact
algebraic identity, $\mu^{-1}=\operatorname{adj}(\varepsilon)$
(\autoref{thm:adjugate}), which makes the impedance product isotropic in every
principal direction and thereby forces the photon dispersion to a doubled
transverse root. Gauge-sector vacuum birefringence about $(1,1,-1)$ therefore
\emph{cancels}, and does so independently of the undetermined coupling
normalizations and of the specific lattice geometry. The companion Paper~IV
confirms the verdict from an independent implementation---matching blocks and a
null split to $\sim10^{-15}$---so the two papers stand together on the same result
by two routes.

What changed by the end is that a question Paper~IV could only pose is now
answered: the framework \emph{predicts the absence} of gauge-sector vacuum
birefringence---a positive null, consistent with polarimetry rather than in
tension with it. Two things are deliberately left open. The shared photon speed
carries a factor-$\sim$2 \emph{common-mode} anisotropy about the optical axis (the
axis is the one cube body-diagonal the hop set omits, so the lattice is trigonal
rather than cubic); this is not birefringence, and whether the physical vacuum
$O_h$-restores to an isotropic speed---rather than realizing a single trigonal
domain---is a genuine open question we do not settle here. Separately, the same
induced action admits a momentum-space route (the physical-band $\mathrm{Tr}\ln T$
eigenphase sum): we fix its mode-filling prescription (the physical band) and
diagnose why a naive band sum diverges (an intraband pole on the exactly-flat
optical axis, Ward-cancelled), but the Ward-safe single-object extraction of the
anisotropic \emph{tensor} is left to a possible follow-on; here the tensor
structure is taken from the plaquette geometry and only the scalar magnitude from
the loop. Neither open item bears on the birefringence verdict.

The natural next steps are those two threads---an isotropy-restoration analysis of
the common-mode speed, and the orbital-susceptibility assembly of the dynamical
tensor---together with the one quantity Paper~I left open and this paper inherits,
the absolute one-loop coupling $1/g^2$, which sets the induced-action magnitude but
not the birefringence verdict.
